% Generated draft: requires mathematical and source review.
\documentclass[11pt]{article}
\usepackage[margin=25mm]{geometry}
\usepackage{fontspec}
\setmainfont{Latin Modern Roman}
\usepackage{amsmath,amssymb,mathtools,unicode-math}
\setmathfont{Latin Modern Math}
\usepackage{enumitem,needspace,xurl,hyperref}
\setlength{\parindent}{0pt}
\setlength{\parskip}{5pt}
\setlength{\emergencystretch}{3em}
\newcommand{\N}{\mathbb{N}}
\newcommand{\Z}{\mathbb{Z}}
\newcommand{\Q}{\mathbb{Q}}
\newcommand{\R}{\mathbb{R}}
\newcommand{\C}{\mathbb{C}}
\newcommand{\F}{\mathbb{F}}
\newcommand{\A}{\mathbb{A}}
\newcommand{\PP}{\mathbb{P}}
\newcommand{\E}{\mathbb{E}}
\newcommand{\eps}{\varepsilon}
\newcommand{\dd}{\,\mathrm{d}}
\begin{document}
\section*{UnsolvedMath ID: 4000018}
\textbf{Source problem code:} catalogue entry
\subsection{Definitions and mathematical statement}
Let \((X,d)\) carry Markov evolution laws \(m_x^{*t}\) starting at \(x\), at the admissible times \(t\) of the evolution. On the usual interpretation, \(T_1\) denotes the Wasserstein transport distance \(W_1\) for \(d\). The displayed strengthened estimate is
\[
T_1(m_x^{*t},m_{x'}^{*t'})\le
 e^{-\kappa\min(t,t')}d(x,x')+
 C\frac{(\sqrt t-\sqrt{t'})^2}{2d(x,x')},
\]
for all distinct \(x,x'\) and admissible nonnegative times \(t,t'\), with \(\kappa>0\) and finite \(C>0\). Distinctness is required because of the denominator. The source associates to this hypothesis the diameter bound
\[
\operatorname{diam}(X)=\sup_{x,x'\in X}d(x,x')
\le\pi\sqrt{\frac{C}{2\kappa}}.
\]
The question concerns the meaning of \(C\), rather than a new proof of that bound: does it reflect an intrinsic dimension, especially the parameter \(n\) in Bakry--Émery \(CD(K,n)\)?

For comparison, in a setting with a suitable Markov generator \(L\), define
\[
\Gamma(f,g)=\tfrac12\bigl(L(fg)-fLg-gLf\bigr),\qquad
\Gamma(f)=\Gamma(f,f),\qquad
\Gamma_2(f)=\tfrac12L\Gamma(f)-\Gamma(f,Lf).
\]
The finite-dimensional curvature-dimension inequality is
\[
\Gamma_2(f)\ge K\Gamma(f)+\frac{(Lf)^2}{n}
\]
pointwise for all admissible test functions, with a specified finite dimension parameter \(n>0\). Comparing these parameters requires the same generator/time normalization and a precise class of spaces; neither \(K=\kappa\) nor \(C=n\) is asserted by the source.

Moreover, any larger \(C\) also satisfies the transport estimate when one \(C\) does. An intrinsic interpretation would therefore have to specify an optimal coefficient, or another normalization selecting \(C\). The source asks for a relation without defining this selection or a proposed formula. The supplied research summary calls the diameter bound established and the specific dimension identification unresolved, while distinguishing known Bakry--Émery and discrete comparison results.
\subsection{Short English statement}
Does the coefficient \(C\) in the stated two-time transport bound represent a dimension, and how does it relate to the \(n\) in \(CD(K,n)\)? The transport and generator normalizations and the choice of coefficient must be fixed.
\subsection{Sources}
\begin{itemize}
\item Local UnsolvedMath registry, permanent ID \(4000018\), exact source code \texttt{catalogue entry}, ``L2 Bonnet--Myers and dimension.''
\item Yann Ollivier, \emph{Discrete Ricci curvature: Open problems} (2008), Problem R; \url{https://www.yann-ollivier.org/rech/publs/problems_curvmarkov.pdf}. Recorded access 2026-07-29; pdf-text extraction.
\item Supplied research summary attributes the displayed Bonnet--Myers-type bound to Ollivier's article, analogous Bakry--Émery bounds to Bakry--Qian, and later discrete Laplacian-comparison bounds to Münch. Complete references and exact normalization conventions are not supplied.
\end{itemize}
\end{document}
